\part{Forms and currents on Berkovich spaces}
\label{part:cld}

This part constructs the local differential theory associated to
Chambert--Loir and Ducros.  Its organization is not that of the original
source.  We begin with calculations on Euclidean polyhedra, introduce
tropical spaces only when those calculations must be made invariant, and
then transport the resulting calculus to Berkovich spaces.  Currents and psh
functions occur before the most elaborate metric formalism, so that each
abstraction is introduced together with the operation it controls.

The final chapters explain the extensions that arose from the limits of the
original smooth-form category.  Gubler--K\"unnemann delta-forms incorporate
polyhedral curvature and tropical intersections; Gubler--Jell--Rabinoff
harmonic tropicalizations enlarge the supply of charts and improve the
cohomological behaviour of forms; and the Gubler--Rabinoff curve-test theory
provides a local language for pluriharmonic, classically psh, and semipositive
piecewise linear functions.  These constructions are compared with CLD forms
and with global model metrics only on their proved common domains.  Similar
terminology is never used as a substitute for a comparison theorem.

The current pluriharmonic comparison includes arbitrary singular points.
By \Cref{thm:supplied-p061-singular-pluriharmonic}, every classically
pluriharmonic function on a good, separated, strict, boundaryless analytic
space over a complete nontrivially valued field is locally
\(\mathbf R\)-piecewise linear. This supplied theorem retires P06.1
in that scope. The earlier quasi-smooth, finite curve-product and toric
arguments remain useful explicit proofs, with their separate boundary
and separation hypotheses recorded in \Cref{app:ch18-frontiers}.
For an arbitrary continuous function with zero CLD Hessian current,
classical pluriharmonicity remains the residual problem of
\Cref{prob:frontier-pluriharmonic}; it is a premise of the supplied
theorem.

\paragraph{Prerequisites and support.}
\label{front:part2-prerequisites-and-support}
The reader needs multivariable calculus and the elementary theory of
differential forms.  The analytic input is supplied by
\Cref{chap:valuations,chap:berkovich,chap:metrics}; the necessary polyhedral
and lattice algebra is developed in the part itself and summarized in
\Cref{app:tropical-algebra}.  Readers interested in currents and weak limits
should also know the Radon-measure material of
\Cref{chap:topological-tools}.  No prior tropical intersection theory is
assumed.

\paragraph{Conceptual clusters.}
\label{front:part2-conceptual-clusters}
The Euclidean and tropical cluster is
\Cref{chap:superforms,chap:tropical-spaces,chap:tropical-integration}:
superforms supply the bigrading and positivity cones, calibrated tropical
spaces supply the domains of integration, and Stokes' formula fixes the
boundary signs.  Analytic tropicalization in
\Cref{chap:analytic-tropicalization} provides the passage to Berkovich
spaces.  The differential cluster
\Cref{chap:cld-forms,chap:cld-currents} constructs forms, test-form
topologies, currents, positivity, and cycle currents.  The potential-theory
cluster
\Cref{chap:cld-poincare-lelong,chap:cld-bedford-taylor,chap:cld-metrized-bundles}
develops Poincar\'e--Lelong, local psh functions, Bedford--Taylor products,
and curvature currents.  Finally,
\Cref{chap:delta-forms,chap:cld-examples} compares delta-forms, harmonic
tropicalizations, curve tests, and the model calculations on curves, tori,
and abelian varieties.

\paragraph{Reading routes.}
\label{front:part2-reading-routes}
The \emph{short route} to forms and currents is
\Cref{chap:superforms,chap:tropical-spaces,chap:tropical-integration,chap:analytic-tropicalization,chap:cld-forms,chap:cld-currents}.
The \emph{standard route} adds
\Cref{chap:cld-poincare-lelong,chap:cld-bedford-taylor,chap:cld-metrized-bundles}
and then uses \Cref{chap:cld-examples} as a consistency check.  The
\emph{reference route} reads the part in order, including
\Cref{chap:delta-forms}; it is the route needed for the comparison and
research boundary in
\Cref{chap:master-dictionary,chap:compare-ma,chap:verified-extensions}.  On a
first pass, the reader may postpone the full calibration formalism, but not
the fact that every integral and current carries a lattice and multiplicity
normalization.

\paragraph{Main outcomes and interfaces.}
\label{front:part2-main-outcomes-and-interfaces}
By the end of the part, the reader should be able to compute
\(\dprime\), \(\ddprime\), and \(\dprimeddprime\) on a tropical chart; integrate
top-degree forms with the correct calibration; pass between smooth forms and
currents; apply the local Poincar\'e--Lelong and Bedford--Taylor formalisms on
their stated domains; and recognize the curvature current of a metrized line
bundle.  The model-metric interface with BFJ is used in
\Cref{chap:compare-ma}.  The Gubler--Rabinoff interface concerns local curve
tests and harmonic charts, whereas the BJ theory of \Cref{part:bj} is global
and valuation-theoretic, with a scaling action; these are distinct comparison
problems.

The approximation-generated CLD psh hierarchy collapses at its first
decreasing-limit stage on every good, topologically separated Berkovich
analytic space over a complete non-Archimedean field, including a trivially
valued field, by \Cref{thm:supplied-p071a-analytic-first-stage}.
Every germ has a decreasing ordinary smooth psh sequence on one fixed
ordinary open neighbourhood, and its limit may be \(-\infty\).
The toric and star calculations remain concrete constructions of this
approximation. The exact CLD star \(X_B\) of
\cite[\S6.7.3]{CLD25} admits the construction for every nonempty branch set
\(B\), by \Cref{thm:p071-infinite-star-collapse-candidate}.
The analytic affine structure is essential: the compact affine
\(G\)-tropical example in \Cref{thm:supplied-p071b-second-stage-star}
has a second-stage germ outside the first stage. Thus a statement for
general affine tropical spaces requires its own hypotheses.

\paragraph{Normalization and domain warnings.}
\label{front:part2-normalization-and-domain-warnings}
The max-plus tropical convention and minus-logarithm analytic coordinates
are reconciled in \Cref{chap:tropical-spaces,app:normalizations}.
Stokes and pushforward formulae retain boundary orientations, primed and
double-primed contractions, calibrations, and multiplicities. Smooth CLD
positivity, current positivity, semipositive PL metrics, classical
curve-test psh functions, and the transfinite CLD-v2 sheaf agree only under
the continuity, PL, dimension, and boundary hypotheses of a cited
comparison theorem.
